\section{Data and preprocessing}
\label{data}

\subsection{Quality control and related-ad grouping}

\textbf{Raw data.} EMSCAD has 17,880 rows, 18 columns and 866 fraudulent ads~\cite{vidros2017}. Each row contains five
free-text fields, categorical metadata and binary source flags. URLs, emails and phone numbers are
represented by hashed placeholders. The audit also records missing values, duplicate content and
class or category support before constructing features.

\textbf{Cleaning.} Spreadsheet errors and other explicit junk markers become missing values.
Original email and phone identities are extracted for grouping before URL, email and phone
placeholders are genericized in model text. Deduplication then compares all posting-content columns,
excluding the row identifier and label, and removes 283 repeated rows after these repairs. No exact
duplicate block has conflicting labels. The cleaned frame contains 17,597 ads, of which 856
(4.864\%) are fraudulent.

\textbf{Related-ad groups.} Rows are connected when they share an exact whitespace-normalized,
lowercase description, a nonempty company profile normalized in the same way, or an original email
or phone identity. Missing descriptions share an empty-description key and are therefore linked,
whereas empty company profiles receive separate identities and are not linked to one another.
Connections chain transitively, producing 4,482 groups. URLs are excluded because shared websites
create large heterogeneous components; consequently, URL identity is outside the split guarantee.

\subsection{Exploratory findings}

\textbf{Class and group structure.} Figure~\ref{fig:eda-overview} shows the cleaned class balance,
group-size skew and text-length distributions. The largest group contains 723 ads, illustrating why
individual rows are not independent sampling units. Considering only nonempty fields, median
description length is 1,011 characters for real ads and 813 for fraudulent ads. Requirement-length
medians are 566 and 352 characters, also among nonempty fields. These summaries describe EMSCAD and
do not establish causal fraud signals.

\begin{figure}[htbp]
\centering
\includegraphics[width=0.96\linewidth]{eda_overview.pdf}
\caption{Cleaned-data overview: (a) class balance; (b) ads by exact related-group size;
(c, d) within-class text-length histograms for nonempty fields on logarithmic character scales,
normalized as \% of nonempty ads in class per bin.}
\label{fig:eda-overview}
\end{figure}

\textbf{Presence signals.} A company profile is present in 83.9\% of real ads but only 32.2\% of
fraudulent ads. Company logos occur in 81.8\% versus 32.9\%, and email placeholders in 6.8\% versus
20.1\%. Figure~\ref{fig:eda-flags} motivates a small set of length and presence features.
Requirements, benefits and URL presence differ little between classes, so these three exploratory
flags are plotted but excluded from the model feature set.

\begin{figure}[htbp]
\centering
\includegraphics[width=0.76\linewidth]{eda_flags.pdf}
\caption{Candidate presence features by class on the 17,597 cleaned ads. Requirements, benefits
and URL presence are exploratory candidates; the other seven flags enter structured models.}
\label{fig:eda-flags}
\end{figure}

\textbf{Categorical associations.} Figure~\ref{fig:eda-categories} displays the two highest and two
lowest fraud rates among categories with at least 100 ads. Category support varies, and extreme
rates can reflect collection procedures or repeated posters. The plot motivates categorical
encoding, not target-based selection of individual categories.

\begin{figure}[htbp]
\centering
\includegraphics[width=0.98\linewidth]{eda_categories.pdf}
\caption{Highest and lowest category fraud rates with at least 100 cleaned ads per category.
Missing categories are labelled explicitly, and the dashed line marks the overall fraud rate.}
\label{fig:eda-categories}
\end{figure}

\subsection{Model inputs}

\textbf{Shared text.} Title, company profile, description, requirements and benefits are
concatenated in that order, with missing fields replaced by empty strings. The classical TF-IDF
configuration uses English stop words, Unicode accent stripping, sublinear term frequency, word
unigrams and bigrams, minimum document frequency three, and at most 20,000 features.

\textbf{Structured blocks.} Table~\ref{tab:features} lists the text and structured inputs for the
classical models and the dense structured block for late fusion. Nominal missing values receive an
explicit \texttt{missing} category before one-hot encoding. Classical numeric features are
standardized without centering the sparse matrix, and binary features pass through unchanged. The
dense fusion preprocessor has the same feature definitions but partition-specific imputation and
scaling, detailed in Section~\ref{method}.

\begin{table}[htbp]
\centering
\small
\begin{tabular}{l>{\raggedright\arraybackslash}p{0.61\linewidth}r}
\toprule
\textbf{Block} & \textbf{Inputs and transformation} & \textbf{Count} \\
\midrule
Text & Word TF-IDF, unigrams and bigrams & up to 20,000 \\
Nominal & Employment type, experience, education, industry, function; one-hot & 5 columns \\
Numeric & Description length, requirements length; standardized & 2 \\
Binary & Telecommuting, logo, questions, profile, salary, email, phone & 7 \\
\bottomrule
\end{tabular}
\caption{Structured feature blocks. Expanded category counts are learned within each training
partition or cross-validation fold.}
\label{tab:features}
\end{table}

\subsection{Saved grouped partitions}

\textbf{Procedure.} For each seed from 0 through 4, every related-ad group is assigned at random to
train, validation or test with probabilities 0.6, 0.2 and 0.2. The draw is repeated until each
partition holds within two percentage points of its target share of ads and has a fraud rate within
one percentage point of the overall rate. Scikit-learn's \texttt{StratifiedGroupKFold} is not used
because it places large groups in a fixed order: it would put the 723-ad group in the test partition
for every seed. The five partitions contain 10,314--10,818 training, 3,190--3,784 validation and
3,199--3,719 test ads, with 137--187 fraud ads per test partition. The group containing 101 fraud ads
always falls in training, since placing it in validation or test would break the fraud-rate tolerance.
Saved indices include a content-and-label fingerprint so that indices from a different cleaned frame
are rejected.

\textbf{Guarantee and boundary.} Checks confirm complete row coverage and no exact related-ad group
shared by train, validation or test within a partition. The guard does not remove paraphrases,
near-duplicate templates, source dependence or temporal dependence. Across different seeds, train,
validation and test rows can overlap, so the partitions are repeated dependent evaluations rather
than independent replicates. Preprocessors and estimators are fitted on training data, although
full-data EDA informed the initial feature design.

\textbf{Training-only EDA check.} Because the EDA used all cleaned ads, its feature decisions were
recomputed on each saved training partition alone. The four kept derived flags (company profile,
salary, email, phone) differ between classes by 0.066--0.477 on every training partition, and
description and requirements medians are lower for fraud in every one, so these decisions do not
depend on validation or test rows. The three excluded flags are borderline: their class differences
reach 0.060--0.082 on three partitions, so a training-only EDA could have kept one of them there.
